\section{Recitation 4: Pattern Matching in the Interpreter}
\label{sec:rec4}

In \Cref{sec:rec3} we \emph{used} Racket's \code{match} to write \code{factorial} and \code{fibonacci}. This recitation \emph{implements} a small version of it: the calculator with \code{let} gets a \code{match} form, and each part of the interpreter is extended in turn: the types, \code{parse}, \code{subst}, and \code{eval}. The file starts with the \code{gensym} import of \Cref{sec:capavoid}, since its \code{let} uses capture-avoiding lazy substitution.

\subsection{The New Form}

A pattern-matching expression looks like
\begin{lstlisting}[style=repl]
(match e (n1 e1) (n2 e2) ... (_ en))
\end{lstlisting}
where each \code{ni} is a number literal and \code{\_} is the wildcard. To evaluate it:
\begin{enumerate}
  \item evaluate \code{e} to a number, once;
  \item try the clauses top to bottom: \code{ni} matches if it equals that number, and \code{\_} matches anything;
  \item the \emph{first} matching clause wins, and only its expression is evaluated;
  \item if no clause matches, it is an error.
\end{enumerate}
This is Racket's \code{match} restricted to literals and \code{\_}, just enough for \code{factorial}- and \code{fibonacci}-style definitions. There are no pattern variables, so a clause never binds a name; that matters for substitution.

\subsection{Types}

A clause pairs a pattern with an expression. The two kinds of clause carry different data (a number pattern has a number, the wildcard nothing), so each gets its own variant:
\begin{lstlisting}
(define-type PatternE
  [numP (n : number) (e : ArithE)]
  [anyP (e : ArithE)])

(define-type ArithE
  [numC (n : number)]
  [plusC (e1 : ArithE) (e2 : ArithE)]
  [timesC (e1 : ArithE) (e2 : ArithE)]
  [letC (x : symbol) (e1 : ArithE) (e2 : ArithE)]
  [idC (x : symbol)]
  [patternC (e : ArithE) (patterns : (listof PatternE))])
\end{lstlisting}
\code{patternC} holds the expression being matched, the \emph{scrutinee}, and the clauses in order. The two types are \emph{mutually recursive}: an \tyname{ArithE} can contain \tyname{PatternE}s, and each \tyname{PatternE} contains an \tyname{ArithE}. \code{define-type} allows that, in either order.

\subsubsection{Examples} The handout's four examples, in ML-like notation:
\begin{center}
\begin{tabular}{@{}lll@{}}
\toprule
 & program & expected \\
\midrule
1 & \code{match 5 with | 5 -> 0 | 6 -> 1} & 0 \\
2 & \code{match (1 + 2) with | 3 -> 1 | \_ -> x} & 1 \\
3 & \code{match x with | 0 -> z | \_ -> (match y with | 0 -> x | 1 -> y)} & unbound variable \\
4 & \code{let x = 1 in let y = 1 in} (the \code{match} of 3, last clause \code{y + 1}) & 2 \\
\bottomrule
\end{tabular}
\end{center}
\begin{lstlisting}
(define example-2 : ArithE
  (patternC (plusC (numC 1) (numC 2))
            (list (numP 3 (numC 1))
                  (anyP (idC 'x)))))

(define example-4 : ArithE
  (letC 'x (numC 1)
        (letC 'y (numC 1)
              (patternC (idC 'x)
                        (list (numP 0 (idC 'z))
                              (anyP (patternC (idC 'y)
                                              (list (numP 0 (idC 'x))
                                                    (numP 1 (plusC (numC 1) (idC 'y)))))))))))
\end{lstlisting}
Examples 2 and 4 both contain a name that nothing binds (\code{x} and \code{z}), yet neither is an error, because that clause is never chosen. In example~3 the error comes from the scrutinee \code{x}, which is always evaluated.

\subsection{Parsing}

The concrete syntax is \code{(match e clause ...)}, so after the \code{match} keyword, \rkt|(second l)| is the scrutinee and \rkt|(rest (rest l))| is the list of clauses. A helper turns that list of s-expressions into a list of \tyname{PatternE}s. Each clause is a two-element list \code{(pattern expr)}, and the pattern decides which variant to build:
\begin{lstlisting}
(define (parse-patterns (s-pats : (listof s-expression))) : (listof PatternE)
  (cond
    [(empty? s-pats) empty]
    [else
     (let [(s-pat (s-exp->list (first s-pats)))]
       (cond
         [(s-exp-number? (first s-pat))
          (cons (numP (s-exp->number (first s-pat)) (parse (second s-pat)))
                (parse-patterns (rest s-pats)))]
         [(equal? (s-exp->symbol (first s-pat)) '_)
          (list (anyP (parse (second s-pat))))]))]))

(define (parse (s : s-expression)) : ArithE
  (cond
    [(s-exp-number? s) (numC (s-exp->number s))]
    [(s-exp-symbol? s) (idC (s-exp->symbol s))]
    [(s-exp-list? s)
     (let [(l (s-exp->list s))]
       (case (s-exp->symbol (first l))
         [(+) (plusC (parse (second l)) (parse (third l)))]
         [(*) (timesC (parse (second l)) (parse (third l)))]
         [(let) (letC (s-exp->symbol (second l)) (parse (third l)) (parse (fourth l)))]
         [(match) (patternC (parse (second l)) (parse-patterns (rest (rest l))))]))]))
\end{lstlisting}
Three details of the TA solution:
\begin{itemize}
  \item A number clause is \code{cons}ed onto the rest, but a \code{\_} clause returns \rkt|(list (anyP ...))|: parsing \emph{stops} at the first wildcard. Any clause after it could never be reached, so dropping it does not change what the program means: \rkt|(parse '(match 1 (_ 0) (2 3)))| has one clause.
  \item The pattern is \emph{not} parsed as an expression; it must be a number literal or \code{\_}. Anything else is an error, though not one this function raises itself: \code{(match 3 ((+ 1 2) 0))} fails in \code{s-exp->symbol} (\code{not a symbol}), and \code{(match 3 (n 0))} falls off the end of the inner \code{cond} (\code{no matching clause}).
  \item \code{parse-patterns} and \code{parse} call each other, just as the two types refer to each other. \code{match} is a symbol in \emph{our} language, matched by \code{case} like \code{+} and \code{let}.
\end{itemize}

\subsection{Substitution}

\begin{exercise}[Why substitute into \code{match}?]
Why do we need to define substitution for pattern-matching expressions?
\end{exercise}
\begin{solution}
Because \code{let} is implemented by substitution, and the body of a \code{let} can contain a \code{match}. In example~4, the \code{x} and \code{y} bound by the \code{let}s are used in the scrutinees and in the clause bodies. If \code{subst} stopped at a \code{patternC}, those \code{idC}s would never be replaced and \code{eval} would report them as unbound. (In \code{plai-typed} the question does not really arise: \code{type-case} rejects a \code{subst} that has no \code{patternC} case.)
\end{solution}

Substituting into a clause means substituting into its expression; the pattern itself is a number or \code{\_}, with no names in it:
\begin{lstlisting}
(define (subst-pattern (erep : ArithE) (x : symbol)) : (PatternE -> PatternE)
  (lambda (p)
    (type-case PatternE p
      [numP (n e) (numP n (subst erep x e))]
      [anyP (e) (anyP (subst erep x e))])))
\end{lstlisting}
\code{subst-pattern} is \emph{curried}: it takes \code{erep} and \code{x} and returns a function from one clause to one clause. That is exactly the shape \code{map} wants, so the \code{patternC} case is one line:
\begin{lstlisting}
(define (subst (erep : ArithE) (x : symbol) (e : ArithE)) : ArithE
  (type-case ArithE e
    [idC (y) (if (equal? x y) erep (idC y))]
    [numC (n) (numC n)]
    [plusC (e1 e2) (plusC (subst erep x e1) (subst erep x e2))]
    [timesC (e1 e2) (timesC (subst erep x e1) (subst erep x e2))]
    [patternC (e pats) (patternC (subst erep x e) (map (subst-pattern erep x) pats))]
    [letC (y e1 e2)
          (cond
            [(equal? x y) (letC y (subst erep x e1) e2)]
            [else
             (let [(z (gensym))]
               (letC z (subst erep x e1) (subst erep x (subst (idC z) y e2))))])]))
\end{lstlisting}
The \code{letC} case is lecture~4's capture-avoiding \code{subst-lazy-v2}. Note what the \code{patternC} case does \emph{not} need: no shadowing check and no renaming. Those exist for \code{letC} because a \code{let} binds a name; our patterns bind nothing, so a clause body sees exactly the same names as the \code{match} around it, like the children of \code{plusC}.

\subsection{Evaluation}

The helper \code{get-pat-expr} takes the number the scrutinee evaluated to and the clauses, and returns the \emph{expression} of the first clause that matches:
\begin{lstlisting}
(define (get-pat-expr (num : number) (pats : (listof PatternE))) : ArithE
  (cond
    [(empty? pats) (error 'get-pat-expr "unmatched expression")]
    [else (type-case PatternE (first pats)
            [numP (n e) (cond [(equal? num n) e] [else (get-pat-expr num (rest pats))])]
            [anyP (e) e])]))

(define (eval (e : ArithE)) : number
  (type-case ArithE e
    [numC (n) n]
    [plusC (e1 e2) (+ (eval e1) (eval e2))]
    [timesC (e1 e2) (* (eval e1) (eval e2))]
    [patternC (e pats) (eval (get-pat-expr (eval e) pats))]
    [letC (x e1 e2) (eval (subst e1 x e2))]
    [idC (x) (error 'eval "unbound variable!")]))
\end{lstlisting}
\code{eval} evaluates the scrutinee once, picks the clause, and evaluates only that clause's expression.

\subsubsection{Why return the expression, not its value?} \code{get-pat-expr} could evaluate every clause and return numbers. That would be wrong, not just slow: example~4 has a clause \code{(0 z)}, and evaluating it would raise ``unbound variable'' even though the scrutinee is 1 and the clause is never chosen. Returning the unevaluated \tyname{ArithE} means only the chosen clause is evaluated, just as only one branch of an \code{if} is.

\subsubsection{Tracing example 4} \code{g1} stands for whatever name \code{gensym} returns.
\begin{lstlisting}[style=repl]
(let x 1 (let y 1 (match x (0 z) (_ (match y (0 x) (1 (+ 1 y)))))))
letC x:  subst (numC 1) for x in the body.
         The inner let binds y, not x, so y is renamed to g1 first.
---> (let g1 1 (match 1 (0 z) (_ (match g1 (0 1) (1 (+ 1 g1))))))
letC g1: subst (numC 1) for g1 in the body
---> (match 1 (0 z) (_ (match 1 (0 1) (1 (+ 1 1)))))
patternC: scrutinee = 1; (0 z) fails; _ matches
---> (match 1 (0 1) (1 (+ 1 1)))
patternC: scrutinee = 1; (0 1) fails; (1 ...) matches
---> (+ 1 1)
---> 2
\end{lstlisting}
The clause \code{(0 z)} is skipped without being evaluated, which is why the free \code{z} is harmless. Example~3 fails at its first step: the scrutinee is \rkt|(idC 'x)|, no \code{let} has replaced it, and \code{eval} reports it as unbound.

\subsubsection{Clause order} Because \code{get-pat-expr} walks the list from the front and stops at the first match, order matters, and a clause after \code{\_} is dead code:
\begin{lstlisting}
(test (eval (parse '(match 5 (_ 0) (5 1)))) 0)
(test (eval (parse '(match 5 (5 1) (_ 0)))) 1)
(test (eval (parse '(match 5 (5 1) (5 2)))) 1)
(test (eval (parse '(match (* 2 3) (5 0) (6 1) (_ 2)))) 1)
(test/exn (eval (parse '(match 7 (5 0) (6 1)))) "unmatched expression")
\end{lstlisting}
The last test is the run-time ``no matching clause'' error of \Cref{sec:rec3}, now in our own language. Nothing can check for it ahead of time: whether a clause matches depends on a number known only when the program runs.
